% theme: respiration_and_blood_circulation
\MajorQuestion{呼吸と血液の循環}
\SetRowSpaceQanda

\Instruction{次の問いに答えなさい。}
\QQ{細かく枝分かれした気管支の先にある、小さな袋を\NamedBlank{alveolus}という。}{肺胞}
\QQ{血液の液体成分を\NamedBlank{plasma}という。}{血しょう}
\QQ{心臓に戻ってくる血液が流れる血管を\NamedBlank{vein}という。}{静脈}
\QQ{血液の成分のうち、酸素を運ぶはたらきをするものを何というか。}{赤血球}
\QQ{肺で酸素を取り入れ、二酸化炭素を体外に出すはたらきを何というか。}{呼吸}

\Instruction{\strut}
\QQ{\RefBlank{alveolus}のまわりを取り巻いている細い血管を\NamedBlank{capillary}という。}{毛細血管}
\QQ{心臓から肺を通り、心臓に戻る経路を何というか。}{肺循環}
\QQ{赤血球に含まれる、酸素と結びつく赤色の物質を何というか。}{ヘモグロビン}
\QQ{\RefBlank{vein}にあり、血液の逆流を防ぐつくりを何というか。}{弁}
\QQ{心臓の上側にあり、全身や肺から血液が戻ってくる部屋を何というか。}{心房}

\Instruction{\strut}
\QQ{多数の\RefBlank{alveolus}があることで、肺は空気と接する何が大きくなるか。}{表面積}
\QQ{血液の成分のうち、体内に入った病原体を分解するものを何というか。}{白血球}
\QQ{心臓から送り出される血液が流れる血管を\NamedBlank{artery}という。}{動脈}
\QQ{\RefBlank{plasma}が\RefBlank{capillary}からしみ出し、細胞と細胞の間を満たしている液体を何というか。}{組織液}
\QQ{肺や、肺とつながって呼吸を行う器官をまとめて何というか。}{呼吸器官}

\Instruction{\strut}
\QQ{血液の成分のうち、出血したときに血液を固めるはたらきをするものを何というか。}{血小板}
\QQ{酸素を多く含む血液と、二酸化炭素を多く含む血液を、それぞれ何というか。}{動脈血、静脈血}
\QQ{心臓から肺以外の全身を通り、心臓に戻る経路を何というか。}{体循環}
\QQ{細胞が酸素を利用して養分からエネルギーを取り出すはたらきを何というか。}{細胞の呼吸}
\QQ{心臓の下側にあり、全身や肺へ血液を送り出す部屋を何というか。}{心室}
